% GENERATED FILE. Edit canonical lessons, then run npm run generate:books.
% canonical-source-hash: fnv1a64:e05b315efff9722e
% canonical-lessons: HI-C453-anand, HI-C453-utsah, HI-C453-muskan, HI-C453-ullas, HI-C453-prasann

\chapter{Joy and Delight}
\label{ch:hi-joy-and-delight}
\hlchaptermodality{453}

\begin{chapteropening}
\textbf{By the end of this chapter:} \emph{I can say how glad someone was.}

Complete the last lesson of chapter 453: I can say how glad someone was.
\end{chapteropening}

\section[\texorpdfstring{\hi{आनंद} (ānand)}{आनंद (ānand)}]{\hi{आनंद} (ānand) — joy; bliss}

\label{lesson:HI-C453-anand}



\begin{quote}
Before the new one: say the Hindi for to feel, then the Hindi for to well up.
\end{quote}

\subsection*{You'll want to know: \hi{आनंद}}
\textbf{\hi{आनंद}} — \emph{ānand} — \textquotedblleft{}joy; bliss\textquotedblright{}.

\begin{quote}
\hi{बच्चों} \hi{के} \hi{साथ} \hi{खेलकर} \hi{दादा} \hi{को} \hi{बहुत} \hi{आनंद} \hi{आया।}
\end{quote}

\begin{quote}
\emph{baccoṁ ke sāth khelkar dādā ko bahut ānand āyā.}
\end{quote}

\begin{quote}
Grandfather had a wonderful time playing with the children.
\end{quote}

Masculine.

\begin{cousinweb}
\textbf{\hi{आनंद}} (\emph{ānand}) is a tatsama, borrowed from Sanskrit \emph{ānanda}, \textquotedblleft{}happiness, joy\textquotedblright{} (Monier-Williams). \emph{ānand ānā}, \textquotedblleft{}joy comes\textquotedblright{}, is how Hindi says someone enjoyed a thing.
\end{cousinweb}

\begin{grammarlens}[title={What you've built}]
One more word for describing what you lived and felt.
\end{grammarlens}

\subsection*{Your turn}
\begin{itemize}
\raggedright
  \item \emph{Say it:} \emph{ānand}
  \item \emph{Say it:} \emph{ānand}, once more
  \item \emph{Say it:} \emph{baccoṁ ke sāth khelkar dādā ko bahut ānand āyā.} — Grandfather had a wonderful time playing with the children.
  \item \emph{Recall:} say the Hindi for to feel, then the Hindi for to well up, then say \emph{ānand} again
\end{itemize}

\begin{tcolorbox}[breakable,colback=teal!4,colframe=teal!35!black,title={Before you move on}]
What does \hi{आनंद} mean? (\textquotedblleft{}Joy; bliss\textquotedblright{}.) Say it once more.
\end{tcolorbox}
\section[\texorpdfstring{\hi{उत्साह} (utsāh)}{उत्साह (utsāh)}]{\hi{उत्साह} (utsāh) — zeal; enthusiasm}

\label{lesson:HI-C453-utsah}



\begin{quote}
Before the new one: say the Hindi for to well up, then the Hindi for joy.
\end{quote}

\subsection*{You'll want to know: \hi{उत्साह}}
\textbf{\hi{उत्साह}} — \emph{utsāh} — \textquotedblleft{}zeal; enthusiasm\textquotedblright{}.

\begin{quote}
\hi{नई} \hi{नौकरी} \hi{के} \hi{लिए} \hi{मीरा} \hi{में} \hi{बहुत} \hi{उत्साह} \hi{है।}
\end{quote}

\begin{quote}
\emph{naī naukrī ke liye Mīrā meṁ bahut utsāh hai.}
\end{quote}

\begin{quote}
Mira is full of enthusiasm for her new job.
\end{quote}

Masculine.

\begin{cousinweb}
\textbf{\hi{उत्साह}} (\emph{utsāh}) is a tatsama, borrowed from Sanskrit \emph{utsāha}, which Monier-Williams glosses as strength of will, effort and energy (Monier-Williams).
\end{cousinweb}

\begin{grammarlens}[title={What you've built}]
One more word for describing what you lived and felt.
\end{grammarlens}

\subsection*{Your turn}
\begin{itemize}
\raggedright
  \item \emph{Say it:} \emph{utsāh}
  \item \emph{Say it:} \emph{utsāh}, once more
  \item \emph{Say it:} \emph{naī naukrī ke liye Mīrā meṁ bahut utsāh hai.} — Mira is full of enthusiasm for her new job.
  \item \emph{Recall:} say the Hindi for to well up, then the Hindi for joy, then say \emph{utsāh} again
\end{itemize}

\begin{tcolorbox}[breakable,colback=teal!4,colframe=teal!35!black,title={Before you move on}]
What does \hi{उत्साह} mean? (\textquotedblleft{}Zeal; enthusiasm\textquotedblright{}.) Say it once more.
\end{tcolorbox}
\section[\texorpdfstring{\hi{मुस्कान} (muskān)}{मुस्कान (muskān)}]{\hi{मुस्कान} (muskān) — a smile}

\label{lesson:HI-C453-muskan}



\begin{quote}
Before the new one: say the Hindi for joy, then the Hindi for zeal.
\end{quote}

\subsection*{You'll want to know: \hi{मुस्कान}}
\textbf{\hi{मुस्कान}} — \emph{muskān} — \textquotedblleft{}a smile\textquotedblright{}.

\begin{quote}
\hi{दादी} \hi{की} \hi{मुस्कान} \hi{बहुत} \hi{मीठी} \hi{है।}
\end{quote}

\begin{quote}
\emph{dādī kī muskān bahut mīṭhī hai.}
\end{quote}

\begin{quote}
Grandmother has a very sweet smile.
\end{quote}

Feminine.

\begin{cousinweb}
\textbf{\hi{मुस्कान}} (\emph{muskān}) is the noun beside \emph{muskānā}, \textquotedblleft{}to smile\textquotedblright{}, a by-form of \emph{muskurānā}; CDIAL files \emph{muskānā} under \emph{*muss}, a form it reconstructs (CDIAL 10227).
\end{cousinweb}

\begin{grammarlens}[title={What you've built}]
One more word for describing what you lived and felt.
\end{grammarlens}

\subsection*{Your turn}
\begin{itemize}
\raggedright
  \item \emph{Say it:} \emph{muskān}
  \item \emph{Say it:} \emph{muskān}, once more
  \item \emph{Say it:} \emph{dādī kī muskān bahut mīṭhī hai.} — Grandmother has a very sweet smile.
  \item \emph{Recall:} say the Hindi for joy, then the Hindi for zeal, then say \emph{muskān} again
\end{itemize}

\begin{tcolorbox}[breakable,colback=teal!4,colframe=teal!35!black,title={Before you move on}]
What does \hi{मुस्कान} mean? (\textquotedblleft{}A smile\textquotedblright{}.) Say it once more.
\end{tcolorbox}
\section[\texorpdfstring{\hi{उल्लास} (ullās)}{उल्लास (ullās)}]{\hi{उल्लास} (ullās) — elation; high spirits}

\label{lesson:HI-C453-ullas}



\begin{quote}
Before the new one: say the Hindi for zeal, then the Hindi for a smile.
\end{quote}

\subsection*{You'll want to know: \hi{उल्लास}}
\textbf{\hi{उल्लास}} — \emph{ullās} — \textquotedblleft{}elation; high spirits\textquotedblright{}.

\begin{quote}
\hi{त्योहार} \hi{के} \hi{दिन} \hi{पूरे} \hi{गाँव} \hi{में} \hi{उल्लास} \hi{था।}
\end{quote}

\begin{quote}
\emph{tyohār ke din pūre gā̃v meṁ ullās thā.}
\end{quote}

\begin{quote}
On the day of the festival the whole village was in high spirits.
\end{quote}

Masculine.

\begin{cousinweb}
\textbf{\hi{उल्लास}} (\emph{ullās}) is a tatsama, borrowed from Sanskrit \emph{ullāsa}, \textquotedblleft{}light, splendour\textquotedblright{} and also \textquotedblleft{}joy, happiness, merriness\textquotedblright{} (Monier-Williams): joy that shines out.
\end{cousinweb}

\begin{grammarlens}[title={What you've built}]
One more word for describing what you lived and felt.
\end{grammarlens}

\subsection*{Your turn}
\begin{itemize}
\raggedright
  \item \emph{Say it:} \emph{ullās}
  \item \emph{Say it:} \emph{ullās}, once more
  \item \emph{Say it:} \emph{tyohār ke din pūre gā̃v meṁ ullās thā.} — On the day of the festival the whole village was in high spirits.
  \item \emph{Recall:} say the Hindi for zeal, then the Hindi for a smile, then say \emph{ullās} again
\end{itemize}

\begin{tcolorbox}[breakable,colback=teal!4,colframe=teal!35!black,title={Before you move on}]
What does \hi{उल्लास} mean? (\textquotedblleft{}Elation; high spirits\textquotedblright{}.) Say it once more.
\end{tcolorbox}
\section[\texorpdfstring{\hi{प्रसन्न} (prasann)}{प्रसन्न (prasann)}]{\hi{प्रसन्न} (prasann) — pleased; cheerful}

\label{lesson:HI-C453-prasann}



\begin{quote}
Before the new one: say the Hindi for a smile, then the Hindi for elation.
\end{quote}

\subsection*{You'll want to know: \hi{प्रसन्न}}
\textbf{\hi{प्रसन्न}} — \emph{prasann} — \textquotedblleft{}pleased; cheerful\textquotedblright{}.

\begin{quote}
\hi{अरुण} \hi{का} \hi{काम} \hi{देखकर} \hi{उसके} \hi{पिता} \hi{प्रसन्न} \hi{हुए।}
\end{quote}

\begin{quote}
\emph{Aruṇ kā kām dekhkar uske pitā prasann hue.}
\end{quote}

\begin{quote}
Arun's father was pleased when he saw his work.
\end{quote}

\begin{cousinweb}
\textbf{\hi{प्रसन्न}} (\emph{prasann}) is a tatsama, borrowed from Sanskrit \emph{prasanna}, \textquotedblleft{}clear, bright\textquotedblright{} and also \textquotedblleft{}soothed, pleased\textquotedblright{} (Monier-Williams). It is a more formal word than \emph{khuś}.
\end{cousinweb}

\begin{grammarlens}[title={What you've built}]
One more word for describing what you lived and felt.
\end{grammarlens}

\subsection*{Your turn}
\begin{itemize}
\raggedright
  \item \emph{Say it:} \emph{prasann}
  \item \emph{Say it:} \emph{prasann}, once more
  \item \emph{Say it:} \emph{Aruṇ kā kām dekhkar uske pitā prasann hue.} — Arun's father was pleased when he saw his work.
  \item \emph{Recall:} say the Hindi for a smile, then the Hindi for elation, then say \emph{prasann} again
\end{itemize}

\begin{tcolorbox}[breakable,colback=teal!4,colframe=teal!35!black,title={Before you move on}]
What does \hi{प्रसन्न} mean? (\textquotedblleft{}Pleased; cheerful\textquotedblright{}.) Say it once more.
\end{tcolorbox}
